% ID: SEL-D0045-OA-E03
% ESTADO: cotejada. Segunda pasada visual y borrador compilado revisados.
% Fuente: Selectividad/sel_2014_matematicas_aplicadas/EXAMEN -ANDAL-5.pdf
% SHA-256: fe436fc3af8e56d0694316005aedeca2fffefbeb445aa9f60efa0aa6da4cba25
% Página: 1; opción A; ejercicio original 3.

En un Instituto de Educación Secundaria el 40\% de los alumnos juegan al fútbol, el 30\% juegan al baloncesto y el 20\% practican ambos deportes.
\begin{enumerate}[label=\alph*)]
\item \textbf{(1 punto)} \textquestiondown{}Cuál es la probabilidad de que un alumno, elegido al azar, no practique ninguno de los dos deportes?
\item \textbf{(0.75 puntos)} Si un alumno, elegido al azar, juega al fútbol, \textquestiondown{}cuál es la probabilidad de que no juegue al baloncesto?
\item \textbf{(0.75 puntos)} \textquestiondown{}Son independientes los sucesos ``jugar al fútbol'' y ``jugar al baloncesto''?
\end{enumerate}
